%% kapa: exam=2 round=33 year=2022 period=3 subject=compensation_law subject_ko=감정평가 및 보상법규 seq=3 number=3 points=20
%% source: Q-Net 감정평가사 기출문제 (cst00309) artlSeq=5210838 file=raw/1008/5210838/2022년 제33회 감정평가사 2차시험 3교시 문제지(감정평가 및 보상법규).pdf
%% conversion: kordoc 4.24.51 -> tools/gen2.py (auto); review=proofread-gpt-astra; flags=none
\begin{longquestion}{【문제 3】}
감정평가사 甲은 A감정평가법인(이하 ‘A법인’이라 함)에 형식적으로만 적을 두었을 뿐 A법인에서 감정평가사 본연의 업무를 전혀 수행하지 않았고 그 법인의 운영에도 관여하지 않았다. 이에 대해 국토교통부장관은 감정평가관리ㆍ징계위원회의 의결에 따라 사전통지를 거쳐 감정평가사 자격취소처분을 하였다. 처분사유는 ‘甲이 A법인에 소속만 유지할 뿐 실질적으로 감정평가업무에 관여하지 아니하는 방법으로 감정평가사의 자격증을 대여하였다’는 것이었고, 그 법적 근거로 「감정평가 및 감정평가사에 관한 법률」(이하 ‘감정평가법’이라 함) 제27조 제1항, 제39조 제1항 단서 및 제2항 제1호가 제시되었다. 甲은 사전통지서에 기재된 의견제출 기한 내에 청문을 신청하였으나 국토교통부장관은 ‘감정평가법 제13조 제1항 제1호에 따라 감정평가사 자격취소를 하려면 청문을 실시하여야 한다는 규정이 있지만, 명의대여를 이유로 하는 감정평가사 자격취소의 경우에는 청문을 실시하여야 한다는 규정이 없을 뿐 아니라 청문을 실시할 필요도 없다’는 이유로 청문을 실시하지 않았다. 甲에 대한 감정평가사 자격취소처분이 적법한지 설명하시오. (20점) \par
\begin{framedtext}
참조 조문 \par
<감정평가 및 감정평가사에 관한 법률> \par
제13조(자격의 취소) ① 국토교통부장관은 감정평가사가 다음 각 호의 어느 하나에 해당하는 경우에는 그 자격을 취소하여야 한다. \par
1. 부정한 방법으로 감정평가사의 자격을 받은 경우 \par
2. 제39조제2항제1호에 해당하는 징계를 받은 경우 \par
제27조(명의대여 등의 금지) ① 감정평가사 또는 감정평가법인등은 다른 사람에게 자기의 성명 또는 상호를 사용하여 제10조에 따른 업무를 수행하게 하거나 자격증ㆍ등록증 또는 인가증을 양도ㆍ대여하거나 이를 부당하게 행사하여서는 아니 된다. \par
제39조(징계) ① 국토교통부장관은 감정평가사가 다음 각 호의 어느 하나에 해당하는 경우에는 제40조에 따른 감정평가관리ㆍ징계위원회의 의결에 따라 제2항 각 호의 어느 하나에 해당하는 징계를 할 수 있다. 다만, 제2항제1호에 따른 징계는 제11호, 제12호를 위반한 경우 및 제27조를 위반하여 다른 사람에게 자격증ㆍ등록증 또는 인가증을 양도 또는 대여한 경우에만 할 수 있다. \par
9. 제25조, 제26조 또는 제27조를 위반한 경우 \par
② 감정평가사에 대한 징계의 종류는 다음과 같다. \par
1. 자격의 취소 \par
2. 등록의 취소 \par
3. 2년 이하의 업무정지 \par
4. 견책 \par
제45조(청문) 국토교통부장관은 다음 각 호의 어느 하나에 해당하는 처분을 하려는 경우에는 청문을 실시하여야 한다. \par
1. 제13조제1항제1호에 따른 감정평가사 자격의 취소 \par
2. 제32조제1항에 따른 감정평가법인의 설립인가 취소 \par
\end{framedtext}
\end{longquestion}
